\documentclass[10pt,a4paper]{amsart}
\usepackage[margin=19mm]{geometry}
\usepackage{amsmath,amssymb,mathtools}
\usepackage[scr=rsfs,cal=euler]{mathalfa}
\usepackage{xfrac}
\usepackage[unicode,hidelinks,bookmarks=false]{hyperref}
\newcommand{\ie}{i.e.\ }
\newcommand{\cf}{cf.\ }
\newcommand{\resp}{respectively\ }
\newcommand{\Subsection}{\subsection}
\providecommand{\nobreakdash}{\nobreak\hbox{-}}
\setlength{\emergencystretch}{2em}
\multlinegap0pt
\usepackage{amssymb}
\newtheorem{proposition}{Proposition}
\title{A Specialized Simplex Algorithm for Budget-Constrained Total Variation-Regularized Problems}
\author{Dominic Yang}
\date{Source: arXiv:2507.13493v3 (CC BY 4.0)}
\begin{document}
\maketitle

\noindent\emph{Source and licence.} A Specialized Simplex Algorithm for Budget-Constrained Total Variation-Regularized Problems. Version arXiv:2507.13493v3 (CC BY 4.0), distributed by arXiv under the Creative Commons Attribution 4.0 International licence, http://creativecommons.org/licenses/by/4.0/, as recorded in the arXiv licence metadata for this exact version and shown on the abstract page https://arxiv.org/abs/2507.13493v3. Full licence notice: \texttt{licenses/arxiv-2507.13493-CC-BY-4.0.txt}.\par\smallskip

\noindent\textit{Editorial scope.} Counted unit: the article's proposition prop:tv-linear-basic (source0.tex lines 86--94). The article's complete original proof of it is delivered with it (source0.tex lines 95--110, one \texttt{proof} environment of 16 lines). No mathematics is written, repaired or re-derived: the statement and the proof are the article's own lines, in the article's own notation, and no retained line is substituted or re-flowed.  Retained uncounted as the source-local context the proof needs: lines 48--55.  Nothing was omitted from any retained span, so the package carries no omission record.  No retained line contains a citation, so the wrapper carries no bibliography and no bibliography entry.  No retained line contains a figure environment, an image-inclusion command or drawing code, so no image is delivered and the package declares no asset dependency.  Editorial formatting only, in addition: an AMS \texttt{amsart} preamble that restates the article's own declarations and macros used by the retained lines, namely the environments \texttt{proposition} on separate counters, together with the standard packages those lines need.  The retained environments print in this document as Proposition 1 in order of appearance, whereas the article prints the same items with the numbering of its own full document.  The complete native source of the article as distributed in the arXiv e-print for this exact version is bundled unchanged as \texttt{sources/181777/source0.tex}.
\begin{equation}\label{eq:TV-Linear}
\begin{aligned}
    \min_{x} \qquad &\sum_{v \in V} c_v x_v + \sum_{uv \in E}d_{uv}a_{uv} + d_{vu} a_{vu} \\
    \text{s.t.} \qquad & x_u - x_v = a_{uv} - a_{vu},  \quad uv \in E\\
                        &\sum_{v\in V}h_v x_v + s = \Delta,\\
                        & 0 \le x_v \le 1,\quad  a_{uv}, a_{vu}, s \ge 0, \quad v \in V, uv \in E
\end{aligned}
\end{equation}

\begin{proposition} \label{prop:tv-linear-basic}
    Let $x$ be a basic solution of \eqref{eq:TV-Linear}. Let $V_{NB} \subset V$ denote the vertices $v$ where $x_v$ are nonbasic, $E^+$ denote edges $uv \in E$ where $a_{uv}$ is nonbasic, $E^-$ denote edges $uv$ where $a_{vu}$ is nonbasic, $E_{NB} = E^+ \cap E^-$ and $G_{NB} = (V, E_{NB})$ denote the subgraph given by selecting only edges in $E_{NB}$. Let $V= V_1 \cup\dots \cup V_k$ partition the vertices of $G_{NB}$ into connected components. Then the following statements are true:
    \begin{enumerate}
        % \item If $s$ is basic, $|V_{NB}| + |E| + |E| = |V| + |E|$; otherwise $|V_{NB}| + |E^+| + |E^-| + 1 = |V| + |E|$
        \item $E^+ \cup E^- = E$.
        \item If $s$ is basic, $|V_{NB} \cap V_i| = 1$ for $i=1,\ldots,k$; if $s$ is nonbasic, $|V_{NB} \cap V_i| = 1$ for all but one $i$ where $|V_{NB} \cap V_i| = 0$ and $h(S_i) \ne 0$.
        \item Each connected component of $G_{NB}$ is a tree.
    \end{enumerate}
\end{proposition}

\begin{proof}
    The first statement follows from the fact that the columns for $a_{uv}$ and $a_{vu}$ are $(-e_{uv}, 0)$ and $(e_{uv}, 0)$ and are dependent so that both can never be in the basis simultaneously.
    For the second statement, we consider two cases. First assume $s$ is basic. Assume that there is some component which contains no nonbasic vertices, and let $S$ be the set of vertices in this component. Let $N(S)$ denote the set of vertices which are not in $S$ but have a neighbor in $S$. As each column corresponding to a vertex variable $v$ is given by
    \[
        A_{TV,v} = (-\sum_{uv \in \delta^-(v)}e_{uv} + \sum_{vw\in \delta^+(v)}e_{vw}, 1),
    \]
    the sum over all such columns is
    \[
        \sum_{v\in S} A_{TV,v} = (-\sum_{wt \in \delta^-(S)}e_{wt} + \sum_{wt \in \delta^+(S)}e_{wt}, h(S))
    \]
    as for each edge $uv$ between vertices in $S$, the corresponding components cancel out. On the other hand, by assumption, one of the variables $a_{uv}, a_{vu}$ for each edge $uv \in \delta(S)$ is basic, and these columns span $\{(e_{uv},1): uv \in \delta(S)\}$. Therefore, we can find a linear combination producing solely the vector $(0, h(S))$ which is linearly dependent with $(0,1)$, the basis column corresponding to $s$. Hence, we do not have a basis, and each component requires at least one nonbasic vertex.

    When $s$ is nonbasic, the column $(0,1)$ is no longer in the basis. If $h(S) = 0$, the argument above still holds and the columns are not linearly independent, but if $h(S) \ne 0$ the above argument does not work as we cannot eliminate the last component. However, if we have two components $S, S'$ both satisfying $h(S), h(S')\ne 0 $, as above we can produce two linear combinations giving $(0, h(S)$ and $(0, h(S'))$ which are linearly dependent. Hence, when $s$ is nonbasic, all but one component must have at least one nonbasic vertex.

    To prove the third statement, we proceed with a counting argument. Note that as there are $|E| + 1$ equations in \eqref{eq:TV-Linear}, and $|V| + 2|E| + 1$ variables, so that there are $|V| + |E|$ nonbasic variables. As $E^+ \cup E^- = E$, we use up $|E|$ nonbasic variables by picking one of $a_{uv}$ and $a_{vu}$ to be nonbasic for each edge $uv \in E$. This leaves $|V|$ variables which correspond to vertices in $V_{NB}$, edges in $E_{NB}$ and possibly $s$. When $s$ is basic, then $|V_{NB}| + |E_{NB}| = |V|$. If $G_{NB}$ has $k$ components, necessarily  $|E_{NB}| \ge |V|-k$, and each component has a nonbasic vertex so that $|V_{NB}| \ge k.$ This is only possible if $|E_{NB}| = |V|-k$ and $|V_{NB}| = k$ implying each component is a tree, and each tree has one nonbasic vertex. When $s$ is nonbasic, then $|V_{NB}| + |E_{NB}| = |V|-1$ and one tree will not have a nonbasic vertex. 
\end{proof}

\end{document}
